% Generated draft: requires mathematical and source review.
\documentclass[11pt]{article}
\usepackage[margin=25mm]{geometry}
\usepackage{fontspec}
\setmainfont{Latin Modern Roman}
\usepackage{amsmath,amssymb,mathtools,unicode-math}
\setmathfont{Latin Modern Math}
\usepackage{enumitem,needspace,xurl,hyperref}
\setlength{\parindent}{0pt}
\setlength{\parskip}{5pt}
\setlength{\emergencystretch}{3em}
\newcommand{\N}{\mathbb{N}}
\newcommand{\Z}{\mathbb{Z}}
\newcommand{\Q}{\mathbb{Q}}
\newcommand{\R}{\mathbb{R}}
\newcommand{\C}{\mathbb{C}}
\newcommand{\F}{\mathbb{F}}
\newcommand{\A}{\mathbb{A}}
\newcommand{\PP}{\mathbb{P}}
\newcommand{\E}{\mathbb{E}}
\newcommand{\eps}{\varepsilon}
\newcommand{\dd}{\,\mathrm{d}}
\begin{document}
\section*{UnsolvedMath ID: 30005442}
\textbf{Source problem code:} catalogue entry
\subsection{Definitions and mathematical statement}
A \emph{skew brace} is a set \(B\) equipped with two group operations \(+\) and \(\circ\) satisfying
\[
a\circ(b+c)=(a\circ b)-a+(a\circ c)\qquad(a,b,c\in B),
\]
where \(-a\) denotes the inverse in \((B,+)\). Neither group is assumed abelian. A skew brace is \emph{bi-skew} if exchanging the two operations also gives a skew brace. It is \emph{two-sided} if it additionally satisfies
\[
(a+b)\circ c=(a\circ c)-c+(b\circ c)\qquad(a,b,c\in B).
\]
These definitions specify the requested output structures; they do not define the semi-affine input structures or the correspondence mentioned in the source.

\textbf{Source-dependent problem.} Let \(G\) be an arbitrary group. Retain the following four tasks, numbered as in the supplied report passage:
\begin{enumerate}
\item[2.] Find all ``trivially decomposable'' cancellative semi-affine structures on \(G\), with the aim of constructing all cancellative semi-affine structures on \(G\) itself.
\item[3.] Extend the description of cancellative semi-affine structures on groups to the general case.
\item[4.] Find systematic constructions, through groupal semi-affine structures on groups, of skew braces that are not bi-skew braces.
\item[5.] Construct two-sided skew braces through groupal semi-affine structures on groups.
\end{enumerate}
No finiteness or commutativity restriction on \(G\) is stated in the supplied passage.

\textbf{Unresolved qualifications.} The supplied record does not give the data or axioms of a semi-affine structure, the meanings of ``cancellative'', ``groupal'', or ``trivially decomposable'' in this setting, or the correspondence that produces a skew brace. The notion of equivalence for the classification is also unspecified. In task 3, neither the existing description nor the scope of ``the general case'' is supplied. Consequently, these tasks preserve the source targets but cannot yet be expressed as a self-contained classification or construction problem. The cleaned statement summarizes the first two tasks as classification of cancellative semi-affine structures on groups; the supplied evidence does not establish that this summary captures every qualification of the original passage.
\subsection{Short English statement}
For an arbitrary group, classify the source's cancellative semi-affine structures, retaining its intermediate task concerning ``trivially decomposable'' structures and its unspecified extension to a general case. Use groupal semi-affine structures to construct skew braces that fail the bi-skew condition and to construct two-sided skew braces. The semi-affine definitions and construction correspondence are absent from the supplied record, so the formulation remains source-dependent and requires human mathematical review.
\subsection{Sources}
\begin{itemize}
\item Supplied local UnsolvedMath registry record, permanent ID \(30005442\), title ``Semi-Affine Structures and Skew Brace Constructions'', exact source code \texttt{catalogue entry}. The original passage contains Problems 2--5. The record gives proposed year 2023 and does not identify a proposer.
\item \emph{Mini-Workshop: Skew Braces and the Yang--Baxter Equation} (2023), \emph{Oberwolfach Reports}, pp. 537--563, DOI \texttt{10.4171/owr/2023/9}. Supplied report citation and literature-source link labeled ``original Oberwolfach report'': \url{https://doi.org/10.4171/owr/2023/9}. No precise page within this range is supplied for Problems 2--5.
\item Supplied literature source: \emph{Semi-affine structures on groups and semi-braces}, \url{https://www.sciencedirect.com/science/article/abs/pii/S002240492300066X}. The record supplies no authors or further bibliographic details for this reference.
\item Literature assessment dated exactly \texttt{2026-08-22}: the supplied curation labels the problem ``partially solved'' and states that a semi-affine correspondence gives a general categorical description and constructions of cancellative and groupal structures, including skew-brace examples, while the requested complete classifications and systematic constructions remain open. This is supplied dated curation, not an independently verified literature assessment. The background explicitly requests independent verification of time-sensitive status and cited claims.
\item The registry labels the statement \texttt{corrected\_verified} and reports that the official report was checked directly and that the cleaned statement restores definitions and removes extraction artifacts or source-dependent wording. These are supplied provenance claims; this draft has not independently visited or verified either reference, and the supplied record itself does not contain the missing semi-affine definitions.
\end{itemize}
\end{document}
